\section{}

\paragraph{1350年}\eventanchor{YUAN-1350-REGNAL-YEAR}{YUAN-1350-REGNAL-YEAR}　元以元顺帝在位第18年作为诏令、赋役与外交文书的纪年框架。账本以《元史》〈本纪〉及相关志；《元朝秘史》；《元典章》或《通制条格》相关条标明来源路径；该锚点连接编年节点，关于数量、实施和地方后果仍应遵守材料的文体与立场限制。

\paragraph{1351年}\eventanchor{YUAN-1351-EVENT-192}{YUAN-1351-EVENT-192}　红巾军起事，元廷同时发动大规模河工。账本以《元史》〈本纪〉及相关志；《元朝秘史》；《元典章》或《通制条格》相关条标明来源路径；该锚点连接编年节点，关于数量、实施和地方后果仍应遵守材料的文体与立场限制。

\paragraph{1351年}\eventanchor{YUAN-1351-REGNAL-YEAR}{YUAN-1351-REGNAL-YEAR}　元以元顺帝在位第19年作为诏令、赋役与外交文书的纪年框架。账本以《元史》〈本纪〉及相关志；《元朝秘史》；《元典章》或《通制条格》相关条标明来源路径；该锚点连接编年节点，关于数量、实施和地方后果仍应遵守材料的文体与立场限制。

\paragraph{1352年}\eventanchor{YUAN-1352-REGNAL-YEAR}{YUAN-1352-REGNAL-YEAR}　元以元顺帝在位第20年作为诏令、赋役与外交文书的纪年框架。账本以《元史》〈本纪〉及相关志；《元朝秘史》；《元典章》或《通制条格》相关条标明来源路径；该锚点连接编年节点，关于数量、实施和地方后果仍应遵守材料的文体与立场限制。

\paragraph{1353年}\eventanchor{YUAN-1353-REGNAL-YEAR}{YUAN-1353-REGNAL-YEAR}　元以元顺帝在位第21年作为诏令、赋役与外交文书的纪年框架。账本以《元史》〈本纪〉及相关志；《元朝秘史》；《元典章》或《通制条格》相关条标明来源路径；该锚点连接编年节点，关于数量、实施和地方后果仍应遵守材料的文体与立场限制。

\paragraph{1354年}\eventanchor{YUAN-1354-EVENT-193}{YUAN-1354-EVENT-193}　脱脱在镇压起事期间失势被罢。账本以《元史》〈本纪〉及相关志；《元朝秘史》；《元典章》或《通制条格》相关条标明来源路径；该锚点连接编年节点，关于数量、实施和地方后果仍应遵守材料的文体与立场限制。

\paragraph{1354年}\eventanchor{YUAN-1354-REGNAL-YEAR}{YUAN-1354-REGNAL-YEAR}　元以元顺帝在位第22年作为诏令、赋役与外交文书的纪年框架。账本以《元史》〈本纪〉及相关志；《元朝秘史》；《元典章》或《通制条格》相关条标明来源路径；该锚点连接编年节点，关于数量、实施和地方后果仍应遵守材料的文体与立场限制。

\paragraph{1355年}\eventanchor{YUAN-1355-REGNAL-YEAR}{YUAN-1355-REGNAL-YEAR}　元以元顺帝在位第23年作为诏令、赋役与外交文书的纪年框架。账本以《元史》〈本纪〉及相关志；《元朝秘史》；《元典章》或《通制条格》相关条标明来源路径；该锚点连接编年节点，关于数量、实施和地方后果仍应遵守材料的文体与立场限制。

\paragraph{1356年}\eventanchor{YUAN-1356-EVENT-194}{YUAN-1356-EVENT-194}　朱元璋攻取集庆并改名应天。账本以《元史》〈本纪〉及相关志；《元朝秘史》；《元典章》或《通制条格》相关条标明来源路径；该锚点连接编年节点，关于数量、实施和地方后果仍应遵守材料的文体与立场限制。

\paragraph{1356年}\eventanchor{YUAN-1356-REGNAL-YEAR}{YUAN-1356-REGNAL-YEAR}　元以元顺帝在位第24年作为诏令、赋役与外交文书的纪年框架。账本以《元史》〈本纪〉及相关志；《元朝秘史》；《元典章》或《通制条格》相关条标明来源路径；该锚点连接编年节点，关于数量、实施和地方后果仍应遵守材料的文体与立场限制。

\paragraph{1357年}\eventanchor{YUAN-1357-REGNAL-YEAR}{YUAN-1357-REGNAL-YEAR}　元以元顺帝在位第25年作为诏令、赋役与外交文书的纪年框架。账本以《元史》〈本纪〉及相关志；《元朝秘史》；《元典章》或《通制条格》相关条标明来源路径；该锚点连接编年节点，关于数量、实施和地方后果仍应遵守材料的文体与立场限制。

\paragraph{1358年}\eventanchor{YUAN-1358-REGNAL-YEAR}{YUAN-1358-REGNAL-YEAR}　元以元顺帝在位第26年作为诏令、赋役与外交文书的纪年框架。账本以《元史》〈本纪〉及相关志；《元朝秘史》；《元典章》或《通制条格》相关条标明来源路径；该锚点连接编年节点，关于数量、实施和地方后果仍应遵守材料的文体与立场限制。

\paragraph{1359年}\eventanchor{YUAN-1359-REGNAL-YEAR}{YUAN-1359-REGNAL-YEAR}　元以元顺帝在位第27年作为诏令、赋役与外交文书的纪年框架。账本以《元史》〈本纪〉及相关志；《元朝秘史》；《元典章》或《通制条格》相关条标明来源路径；该锚点连接编年节点，关于数量、实施和地方后果仍应遵守材料的文体与立场限制。

\paragraph{1360年}\eventanchor{YUAN-1360-REGNAL-YEAR}{YUAN-1360-REGNAL-YEAR}　元以元顺帝在位第28年作为诏令、赋役与外交文书的纪年框架。账本以《元史》〈本纪〉及相关志；《元朝秘史》；《元典章》或《通制条格》相关条标明来源路径；该锚点连接编年节点，关于数量、实施和地方后果仍应遵守材料的文体与立场限制。

\paragraph{1361年}\eventanchor{YUAN-1361-REGNAL-YEAR}{YUAN-1361-REGNAL-YEAR}　元以元顺帝在位第29年作为诏令、赋役与外交文书的纪年框架。账本以《元史》〈本纪〉及相关志；《元朝秘史》；《元典章》或《通制条格》相关条标明来源路径；该锚点连接编年节点，关于数量、实施和地方后果仍应遵守材料的文体与立场限制。

\paragraph{1362年}\eventanchor{YUAN-1362-REGNAL-YEAR}{YUAN-1362-REGNAL-YEAR}　元以元顺帝在位第30年作为诏令、赋役与外交文书的纪年框架。账本以《元史》〈本纪〉及相关志；《元朝秘史》；《元典章》或《通制条格》相关条标明来源路径；该锚点连接编年节点，关于数量、实施和地方后果仍应遵守材料的文体与立场限制。

\paragraph{1363年}\eventanchor{YUAN-1363-EVENT-195}{YUAN-1363-EVENT-195}　鄱阳湖之战改变江南群雄力量对比。账本以《元史》〈本纪〉及相关志；《元朝秘史》；《元典章》或《通制条格》相关条标明来源路径；该锚点连接编年节点，关于数量、实施和地方后果仍应遵守材料的文体与立场限制。

\paragraph{1363年}\eventanchor{YUAN-1363-REGNAL-YEAR}{YUAN-1363-REGNAL-YEAR}　元以元顺帝在位第31年作为诏令、赋役与外交文书的纪年框架。账本以《元史》〈本纪〉及相关志；《元朝秘史》；《元典章》或《通制条格》相关条标明来源路径；该锚点连接编年节点，关于数量、实施和地方后果仍应遵守材料的文体与立场限制。

\paragraph{1364年}\eventanchor{YUAN-1364-REGNAL-YEAR}{YUAN-1364-REGNAL-YEAR}　元以元顺帝在位第32年作为诏令、赋役与外交文书的纪年框架。账本以《元史》〈本纪〉及相关志；《元朝秘史》；《元典章》或《通制条格》相关条标明来源路径；该锚点连接编年节点，关于数量、实施和地方后果仍应遵守材料的文体与立场限制。

\paragraph{1365年}\eventanchor{YUAN-1365-REGNAL-YEAR}{YUAN-1365-REGNAL-YEAR}　元以元顺帝在位第33年作为诏令、赋役与外交文书的纪年框架。账本以《元史》〈本纪〉及相关志；《元朝秘史》；《元典章》或《通制条格》相关条标明来源路径；该锚点连接编年节点，关于数量、实施和地方后果仍应遵守材料的文体与立场限制。

\paragraph{1366年}\eventanchor{YUAN-1366-REGNAL-YEAR}{YUAN-1366-REGNAL-YEAR}　元以元顺帝在位第34年作为诏令、赋役与外交文书的纪年框架。账本以《元史》〈本纪〉及相关志；《元朝秘史》；《元典章》或《通制条格》相关条标明来源路径；该锚点连接编年节点，关于数量、实施和地方后果仍应遵守材料的文体与立场限制。

\paragraph{1367年}\eventanchor{YUAN-1367-EVENT-196}{YUAN-1367-EVENT-196}　明军发动北伐，元在华北防线动摇。账本以《元史》〈本纪〉及相关志；《元朝秘史》；《元典章》或《通制条格》相关条标明来源路径；该锚点连接编年节点，关于数量、实施和地方后果仍应遵守材料的文体与立场限制。

\paragraph{1367年}\eventanchor{YUAN-1367-REGNAL-YEAR}{YUAN-1367-REGNAL-YEAR}　元以元顺帝在位第35年作为诏令、赋役与外交文书的纪年框架。账本以《元史》〈本纪〉及相关志；《元朝秘史》；《元典章》或《通制条格》相关条标明来源路径；该锚点连接编年节点，关于数量、实施和地方后果仍应遵守材料的文体与立场限制。

\paragraph{1368年}\eventanchor{YUAN-1368-EVENT-197}{YUAN-1368-EVENT-197}　明军进入大都，元顺帝北撤。账本以《元史》〈本纪〉及相关志；《元朝秘史》；《元典章》或《通制条格》相关条标明来源路径；该锚点连接编年节点，关于数量、实施和地方后果仍应遵守材料的文体与立场限制。

\paragraph{1368年}\eventanchor{YUAN-1368-REGNAL-YEAR}{YUAN-1368-REGNAL-YEAR}　元以元顺帝在位第36年作为诏令、赋役与外交文书的纪年框架。账本以《元史》〈本纪〉及相关志；《元朝秘史》；《元典章》或《通制条格》相关条标明来源路径；该锚点连接编年节点，关于数量、实施和地方后果仍应遵守材料的文体与立场限制。
